\sectionheading{5.28 Python Installation ও Setup}

Python-এ Program লেখা ও Execute করার আগে Computer-এ
\textbf{Python Interpreter} Install করতে হয়।

Installation সম্পন্ন হওয়ার পর Python Shell, IDLE অথবা
প্রয়োজনীয় Code Editor/IDE ব্যবহার করে Python Program
লেখা ও চালানো যায়।

%================================================
% Definition
%================================================

\booksubheading{Python Installation কী?}

\begin{definitionbox}

\textbf{\color{darkTeal}সংজ্ঞা}

Computer-এ Python Interpreter এবং প্রয়োজনীয় Supporting Tool
স্থাপন করে Python Program লেখা ও চালানোর উপযোগী পরিবেশ
তৈরি করার প্রক্রিয়াকে \textbf{Python Installation ও Setup} বলা হয়।

\end{definitionbox}

\textbf{সহজভাবে মনে রাখবে—}

\begin{center}

\textbf{Download}
\(\longrightarrow\)
\textbf{Install}
\(\longrightarrow\)
\textbf{PATH Setup}
\(\longrightarrow\)
\textbf{Verify}
\(\longrightarrow\)
\textbf{Run}

\end{center}

%================================================
% Installation Steps
%================================================

\booksubheading{Python Install করার সাধারণ ধাপ}

Windows Computer-এ Python Install করার সাধারণ প্রক্রিয়া—

\begin{enumerate}

\item
Python-এর Official Website থেকে Python 3-এর Installer
Download করতে হবে।

\item
Download করা Installer চালু করতে হবে।

\item
Installation-এর শুরুতে
\textbf{Add python.exe to PATH}
Option থাকলে সেটি নির্বাচন করতে হবে।

\item
এরপর \textbf{Install Now} নির্বাচন করে Installation শুরু করতে হবে।

\item
Installation সম্পন্ন হওয়ার পর Command Prompt বা Terminal
চালু করতে হবে।

\item
Python সঠিকভাবে Install হয়েছে কি না Version দেখে যাচাই করতে হবে।

\item
এরপর IDLE, Python Shell অথবা প্রয়োজনীয় Editor ব্যবহার করে
Program লেখা শুরু করা যাবে।

\end{enumerate}

\replaceableimage[60mm]
{images/python-installer.png}
{Windows-এ Python Installer: Add python.exe to PATH এবং Install Now}

\begin{center}
\footnotesize\textbf{চিত্র: Windows-এ Python Installation-এর প্রাথমিক ধাপ}
\end{center}

%================================================
% PATH
%================================================

\booksubheading{Add python.exe to PATH কেন গুরুত্বপূর্ণ?}

Installation-এর সময় গুরুত্বপূর্ণ একটি Option হলো—

\begin{center}

\textbf{Add python.exe to PATH}

\end{center}

এটি নির্বাচন করলে Operating System Command Prompt বা Terminal
থেকে Python Program চালানোর জন্য Python Interpreter-এর
Location সহজে খুঁজে পায়।

\begin{center}

\textbf{Python PATH-এ যুক্ত}
\(\longrightarrow\)
\textbf{Terminal থেকে Python চালানো সহজ}

\end{center}

\begin{rememberbox}

\textbf{\color{deepOrange}সতর্কতা}

PATH সঠিকভাবে Set না থাকলে Command Prompt-এ
Python Command ব্যবহার করার সময় Error দেখা দিতে পারে।

উদাহরণ—

\begin{codebox}

\begin{english}
\ttfamily
python is not recognized ...
\end{english}

\end{codebox}

প্রয়োজনে পরে Environment Variable-এর PATH-এ Python-এর
Location যোগ করেও এটি ঠিক করা যায়।

\end{rememberbox}

%================================================
% Verify Installation
%================================================

\booksubheading{Python Installation যাচাই করা}

Python Install হওয়ার পর Command Prompt বা Terminal খুলে
নিচের Command ব্যবহার করা যায়—

\begin{codebox}

\begin{english}
\ttfamily
python --version
\end{english}

\end{codebox}

Windows-এর অনেক Setup-এ নিচের Command-ও ব্যবহার করা যায়—

\begin{codebox}

\begin{english}
\ttfamily
py --version
\end{english}

\end{codebox}

Python সঠিকভাবে Install থাকলে Installed Version দেখা যাবে।

উদাহরণ—

\begin{codebox}

\begin{english}
\ttfamily
Python 3.x.x
\end{english}

\end{codebox}

এখানে
\textenglish{\ttfamily 3.x.x}
দ্বারা Python 3-এর Installed Version বোঝানো হয়েছে।

\begin{rememberbox}

\textbf{\color{deepOrange}মনে রাখবে}

Version সঠিকভাবে দেখা যাচ্ছে
\(\longrightarrow\)
Python Installation সফলভাবে হয়েছে।

\end{rememberbox}

%================================================
% Python Shell
%================================================

\booksubheading{Python Shell চালু করা}

Command Prompt বা Terminal-এ লিখতে পারি—

\begin{codebox}

\begin{english}
\ttfamily
python
\end{english}

\end{codebox}

অথবা উপযুক্ত Windows Setup-এ—

\begin{codebox}

\begin{english}
\ttfamily
py
\end{english}

\end{codebox}

এরপর Python Interactive Shell-এর Prompt দেখা যেতে পারে—

\begin{codebox}

\begin{english}
\ttfamily
>>>
\end{english}

\end{codebox}

এই Prompt-এ সরাসরি Python Instruction লেখা যায়।

\begin{examplebox}{উদাহরণ}

\begin{codebox}

\begin{english}
\ttfamily
>>> 5 + 3\\
8
\end{english}

\end{codebox}

এখানে
\textenglish{\ttfamily 5 + 3}
লেখার পর সঙ্গে সঙ্গে
\textenglish{\ttfamily 8}
Result পাওয়া গেছে।

\end{examplebox}

\begin{rememberbox}

\textbf{\color{deepOrange}MCQ তথ্য}

\textenglish{\ttfamily >>>}
\(\longrightarrow\)
Python Interactive Shell-এর Prompt

\end{rememberbox}

%================================================
% IDLE
%================================================

\booksubheading{IDLE চালু করা}

Python-এর Standard Installation-এর সঙ্গে সাধারণত
\textbf{IDLE} ব্যবহার করার সুবিধা পাওয়া যায়।

Windows-এ Start Menu থেকে IDLE খুঁজে চালু করা যায়।

IDLE-এর মাধ্যমে—

\begin{compacttwocolpoints}

\item Python Shell ব্যবহার করা যায়।

\item নতুন Python Program লেখা যায়।

\item Program File Save করা যায়।

\item Program Run করা যায়।

\item Error Message দেখা যায়।

\end{compacttwocolpoints}

\begin{rememberbox}

\textbf{\color{deepOrange}মনে রাখবে}

\textbf{IDLE}
=
Integrated Development and Learning Environment

\end{rememberbox}

%================================================
% Test Program
%================================================

\booksubheading{প্রথম Test Program}

Installation ও Environment সঠিকভাবে কাজ করছে কি না
যাচাই করার জন্য একটি ছোট Program চালানো যায়—

\begin{examplebox}{Test Program}

\begin{codebox}

\begin{english}
\ttfamily
print("Hello, Python!")
\end{english}

\end{codebox}

\textbf{Output:}

\begin{codebox}

\begin{english}
\ttfamily
Hello, Python!
\end{english}

\end{codebox}

এতে বোঝা যায় Python Environment Program Execute করতে পারছে।

\end{examplebox}

%================================================
% File Extension
%================================================

\booksubheading{Python Program File-এর Extension}

Python Program সাধারণত
\textbf{\textenglish{\ttfamily .py}}
Extension-এ Save করা হয়।

উদাহরণ—

\begin{codebox}

\begin{english}
\ttfamily
hello.py\\
sum.py\\
program.py
\end{english}

\end{codebox}

\begin{rememberbox}

\textbf{\color{deepOrange}MCQ তথ্য}

Python Source File Extension
\(\longrightarrow\)
\textbf{\textenglish{\ttfamily .py}}

\end{rememberbox}

%================================================
% Setup
%================================================

\booksubheading{Setup বলতে কী বোঝায়?}

শুধু Python Install করলেই Programming-এর প্রস্তুতি সম্পূর্ণ হয় না।

Programming শুরু করার উপযোগী অবস্থায় প্রয়োজনীয় Environment
প্রস্তুত করাও Setup-এর অংশ।

\begin{tabularx}
{\linewidth}
{
|>{\bfseries\RaggedRight\arraybackslash}p{0.34\linewidth}
|Y|
}

\hline

\rowcolor{softBlue}

\multicolumn{1}{|c|}{\textbf{কাজ}}
&
\multicolumn{1}{c|}{\textbf{উদ্দেশ্য}}

\\

\hline

Python Install
&
Python Interpreter পাওয়া
\\

\hline

PATH Setup
&
Terminal থেকে সহজে Python চালানো
\\

\hline

Version Check
&
Installation সঠিক হয়েছে কি না যাচাই করা
\\

\hline

IDLE / Editor চালু
&
Code লেখার পরিবেশ প্রস্তুত করা
\\

\hline

Test Program Run
&
Python Environment ঠিকভাবে কাজ করছে কি না যাচাই করা
\\

\hline

\end{tabularx}

%================================================
% Common Installation Problems
%================================================

\booksubheading{Installation-এর সাধারণ সমস্যা}

\begin{tabularx}
{\linewidth}
{
|>{\RaggedRight\arraybackslash}p{0.34\linewidth}
|Y|
}

\hline

\rowcolor{softBlue}

\multicolumn{1}{|c|}{\textbf{সমস্যা}}
&
\multicolumn{1}{c|}{\textbf{সম্ভাব্য কারণ/সমাধান}}

\\

\hline

\textenglish{\ttfamily python}
Command কাজ করছে না
&
PATH সঠিকভাবে Set হয়েছে কি না পরীক্ষা করতে হবে।
\\

\hline

Python Version দেখা যাচ্ছে না
&
Installation সম্পূর্ণ হয়েছে কি না পরীক্ষা করতে হবে।
\\

\hline

পুরোনো Version চালু হচ্ছে
&
Computer-এ একাধিক Python Version থাকতে পারে।
\\

\hline

Program Run হচ্ছে না
&
Command, File এবং Environment Setup পরীক্ষা করতে হবে।
\\

\hline

IDLE পাওয়া যাচ্ছে না
&
Python Installation ও IDLE Component পরীক্ষা করতে হবে।
\\

\hline

\end{tabularx}

%================================================
% Python 3
%================================================

\booksubheading{কোন Version ব্যবহার করবে?}

নতুন Program শেখা ও এই বইয়ের উদাহরণের জন্য
\textbf{Python 3} ব্যবহার করা হবে।

এই বইয়ের পরবর্তী সকল Python Program
\textbf{Python 3 Syntax} অনুসরণ করবে।

\begin{rememberbox}

\textbf{\color{deepOrange}দ্রুত ধারাবাহিকতা}

\begin{center}

\textbf{Download}
\(\longrightarrow\)
\textbf{Add to PATH}
\(\longrightarrow\)
\textbf{Install}
\(\longrightarrow\)
\textbf{Version Check}
\(\longrightarrow\)
\textbf{Run}

\end{center}

\end{rememberbox}

%================================================
% Common Mistakes
%================================================

\booksubheading{সাধারণ ভুল}

\begin{compacttwocolpoints}

\item Python Install না করেই Local Computer-এ Program
Execute করার চেষ্টা করা।

\item
\textbf{Add python.exe to PATH}
Option উপেক্ষা করা।

\item Python Installation হয়েছে কি না Version Check না করা।

\item Python Source File ভুল Extension-এ Save করা।

\item Python 2 ও Python 3-এর Syntax মিশিয়ে ব্যবহার করা।

\item Python, IDLE এবং Code Editor-কে একই Software মনে করা।

\item
\textenglish{\ttfamily >>>}
-কে Program-এর অংশ মনে করে File-এর মধ্যে লিখে দেওয়া।

\end{compacttwocolpoints}

%================================================
% Final Reminder
%================================================

\begin{rememberbox}

\textbf{\color{deepOrange}মনে রাখবে}

\begin{itemize}

\item
Python Program চালাতে Computer-এ
\textbf{Python Interpreter}
Install থাকতে হয়।

\item
এই বইয়ের জন্য
\textbf{Python 3}
ব্যবহার করা হবে।

\item
Installation-এর সময়
\textbf{Add python.exe to PATH}
Option গুরুত্বপূর্ণ।

\item
Installation যাচাই করতে—

\begin{center}

\textenglish{\ttfamily python --version}

অথবা

\textenglish{\ttfamily py --version}

\end{center}

ব্যবহার করা যায়।

\item
\textenglish{\ttfamily >>>}
হলো Python Interactive Shell-এর Prompt।

\item
Python Source File-এর Extension হলো
\textbf{\textenglish{\ttfamily .py}}।

\item
IDLE ব্যবহার করে Python Code লেখা, Save এবং Run করা যায়।

\item
দ্রুত মনে রাখবে—

\begin{center}

\textbf{Install Python}
\(\longrightarrow\)
\textbf{PATH ঠিক করো}
\(\longrightarrow\)
\textbf{Version Check করো}
\(\longrightarrow\)
\textbf{Test Program চালাও}

\end{center}

\end{itemize}

\end{rememberbox}

\small

%================================================
% 5.29 Ways of Writing Python Code
%================================================

